\fronttitle{Foreword and introduction}
\label{Intro}

\begin{center}
{\large\bfseries Foreword}
\end{center}

We present here a slightly revised re-edition of the original Seminar, whose
aim and content are indicated in the Introduction. The revision consisted
essentially in the correction of misprints, the addition (in footnotes) of a
few supplementary remarks or references, the present division into three
volumes each provided with a detailed table of contents and an index of
notations, and the addition of a terminological index at the end of
volume~3. Moreover, expos\'e~$\mathrm{VI}_{\mathrm{B}}$ of J.-E.~Bertin has
been partially rewritten by him, notably paragraphs~5 and~10, so that certain
references to this expos\'e differ from the references to the original
expos\'e. The reader will find a list of the expos\'es of the Seminar at the
beginning of the present volume.

Since the publication of the first edition of the present Seminar, the whole
of the \emph{\'El\'ements de G\'eom\'etrie Alg\'ebrique}, Chap.~IV, has
appeared, which renders certain passages of the Seminar unnecessary; we have
sometimes indicated in footnotes the relevant references to EGA~IV which make
it possible to bypass such passages.

For another expos\'e on algebraic groups making systematic use of the
language of schemes, we mention the book of M.~Demazure and P.~Gabriel,
\emph{Groupes alg\'ebriques} (North-Holland \& Masson et Cie). Unlike the
present Seminar, this book assumes no knowledge of algebraic geometry, but
contains all the necessary preliminaries of scheme theory, and its reading
can therefore serve as an introduction to the study of our Seminar. (It
contains, moreover, themes not covered in the Seminar, such as the structure
theory \`a la Dieudonn\'e of commutative affine algebraic groups, in
\loccit Chap.~V.)

\medskip
\begin{flushright}
Bures-sur-Yvette, March~1970
\end{flushright}

\begin{center}
{\large\bfseries Introduction}
\end{center}

\par\addvspace{0.55\baselineskip}
\noindent\textbf{1.}\label{Intro.sec.1}
The aim of the present seminar is twofold.

On the one hand, we aim to give convenient foundations for the theory of
group schemes in general. Expos\'es~I to~IV will give, in this respect, the
indispensable preliminary exercises in schematic and categorical syntax. In
order to obtain a language which ``sticks'' without effort to the geometric
intuition, and to avoid circumlocutions which are unbearable in the long
run, we always identify a prescheme~$X$ over another~$S$ with the functor
\(\Sch^{\circ}/S\to\Ens\) which it represents%
{\renewcommand{\thefootnote}{(*)}\footnote{Such a point of view seems to have
been envisaged for the first time eight or nine years ago, in connection
with the theory of formal groups, by P.~Cartier, who unfortunately did not
take the trouble to make it precise and to systematize it as it
deserved.}},
and it is necessary to give numerous definitions in such a way that they
apply to arbitrary functors, representable or not. Moreover, almost all the
functors which we shall have to use will be ``sheaves'' (for the ``faithfully
flat quasi-compact topology''); expos\'e~IV, which treats groups only incidentally, gives a sketch of the language of ``localization'' and of
sheaves, which likewise proves very convenient in questions of
representability of functors. This expos\'e will above all provide us, for
questions of passage to the quotient, with the most convenient framework for
the sequel.

Expos\'e~V gives some general results on the existence of quotients, taken
up again in expos\'e~$\mathrm{VI}_{\mathrm{A}}$ in the case of the quotient
of an algebraic group over a field (or, more generally, over an artinian
ring) by a subgroup%
{\renewcommand{\thefootnote}{(**)}\footnote{For a more thorough study of
passage to the quotient, notably by reductive groups, see the important
study of D.~Mumford, \emph{Geometric Invariant Theory}, Ergebnisse der
Mathematik, Bd~34, Springer 1965. Let us observe that on one important
point, the terminology of this book does not agree with ours, for over a
field of characteristic \(p > 0\), the groups which Mumford calls
``reductive''\textsuperscript{(1)}, are the smooth groups of multiplicative
type in the sense of the Seminar (\cf Exp.~IX). One may doubtless consider
that the sense in which Mumford takes the word ``reductive'', which loses
its meaning over a base which is not a field, was adopted by him
provisionally and as a sort of makeshift (and this is also more or less what
Mumford explains, for other reasons, from the second paragraph of his
preface!). \textsuperscript{(1)}~\textbf{N.D.E.} see the remarks added at
the end of this Introduction.}}.
This latter expos\'e and expos\'e~$\mathrm{VI}_{\mathrm{B}}$, which follows
it, also contain various elementary results special to algebraic groups over
a field, commonly used in the sequel.

Expos\'e~VII studies certain facts connected with the characteristic of the
base field and develops in particular, with the appropriate generality, the
correspondence between radicial group schemes of height~\(1\) and restricted
Lie \(p\)-algebras.

Finally, expos\'e~XVIII contains the generalization, in the theory of
schemes, of Weil's theorem on the ``birational'' definition of algebraic
groups.

On the other hand, we propose to generalize, to groups over an arbitrary
base prescheme, the Borel--Chevalley structure theory of affine algebraic
groups. It has moreover become apparent, on the occasion of the writing up
of the notes of the seminar, that the affine hypothesis was unnecessary for
many results of the theory. The most complete results are of course obtained
in the cases of ``semi-simple group schemes'', or more generally
``reductive'' ones, with which we shall concern ourselves exclusively from
expos\'e~XIX on. Chevalley himself had already given the construction of the
groups ``of T\^ohoku'' over the ring of integers, a construction which will
be taken up again in the present seminar. The principal uniqueness
theorem%
{\renewcommand{\thefootnote}{(2)}\nde{\label{Intro.fn.2}This refers to corollary~XXIII.5.6,
which follows easily from the uniqueness theorem for split reductive groups
(\cf XXIII, th.~4.1 and cor.~5.2), given that every semi-simple $S$-group is
a ``twisted form'' (for the \'etale topology) of a group ``of T\^ohoku''
(\cf XXII, cor.~2.3).}}
gives a simple characterization of the ``twisted'' variants of these T\^ohoku
groups, over a base prescheme~$S$: these are the groups which are affine and
smooth over~$S$, whose geometric fibers are connected semi-simple groups in
the usual sense%
{\renewcommand{\thefootnote}{(***)}\footnote{This is the essential result of
the thesis of M.~Demazure (\emph{Sch\'emas en groupes r\'eductifs}, Bull.\
Soc.\ Math.\ France \textbf{93} (1965), 369--413).}}\,\textsuperscript{(2)}.

\par\addvspace{0.55\baselineskip}
\noindent\textbf{2.}\label{Intro.sec.2}
As in the case of the known theory over an algebraically closed field, a
crucial role is played by the subtori of the group schemes under
consideration. Thus the preliminary study of tori, and more generally of
``group schemes of multiplicative type'', (both from the intrinsic point of
view and from the point of view of the subgroups of multiplicative type of a
given group), takes a rather large place in this Seminar
(expos\'es~VIII to~XII). Their remarkable rigidity (greater even, in certain
respects, than that of abelian schemes, or of semi-simple group schemes)
makes them very effective working tools for the study of certain more
general groups.

\par\addvspace{0.55\baselineskip}
\noindent\textbf{3.}\label{Intro.sec.3}
From expos\'e~XII on (to the exclusion of expos\'e~XVIII, already mentioned)
we shall commonly use the theory of affine algebraic groups over an
algebraically closed field, which the reader will find in the Chevalley
Seminar~1956, more particularly in expos\'es~IV to~IX of this Seminar. We
shall also use, but to a lesser extent, the later expos\'es of the Chevalley
Seminar, devoted to the structure of semi-simple algebraic groups. Indeed,
we shall take up Chevalley's theory directly in the framework of schemes:
one will see that in this way (even over a base field) the expos\'e gains in
simplicity and in precision.

\par\addvspace{0.55\baselineskip}
\noindent\textbf{4.}\label{Intro.sec.4}
The principal object of the present Seminar is of course to develop
techniques which apply to the study of group schemes over an arbitrary base,
\ie essentially to the study of families of algebraic groups. As such, the
infinitesimal properties of such families, and in particular the case of an
artinian base scheme, play an important role. These properties intervene
even for the study of algebraic groups over a field~$k$, in the case where
the latter is not perfect, so as to be able in particular to apply the
technique of descent in the non-Galois case. Among the new results obtained
in this case, let us mention the existence of maximal tori and of Cartan
subgroups in an arbitrary smooth algebraic group, the rationality of the
variety of maximal tori, and various related results
(Exp.~XIV
{\renewcommand{\thefootnote}{(3)}\nde{in particular, theorems~1.1 and~6.1.}}),
or the correspondence between the ``forms'' of a semi-simple group and the
principal homogeneous bundles under a suitable semi-simple algebraic group
(in general not connected) (Exp.~XXIV, 1.16--1.20). In general, one may say
that the methods required in order to work over a base field which is not
perfect are essentially those used for arbitrary base preschemes, and
thereby fall outside the framework of classical algebraic geometry.

\par\addvspace{0.55\baselineskip}
\noindent\textbf{5.}\label{Intro.sec.5}
It did not seem useful to indicate at the head of the written expos\'es the
date or the dates of the corresponding oral expos\'es of the Seminar. Let us
be content to say that the order of the mimeographed expos\'es (from~I
to~XXVI) corresponds well to the order of the oral expos\'es. Moreover, the
writing of the definitive text is sometimes distinctly later than that of
the oral expos\'e, and often differs from it rather substantially, the
written text being generally more detailed and more complete (such as
Exp.~IV and~$\mathrm{VII}_{\mathrm{B}}$), or even appreciably more general
(such as Exp.~XII or~$\mathrm{VII}_{\mathrm{B}}$) than the oral expos\'e.
Other written expos\'es correspond to no oral expos\'e
($\mathrm{VI}_{\mathrm{B}}$, $\mathrm{VII}_{\mathrm{A}}$, XV, XVI, XVII, and
the essential part of~XXVI), and were written and inserted into the
mimeographed Seminar, either to provide convenient references for various
other expos\'es (this is notably the case of~$\mathrm{VI}_{\mathrm{B}}$), or
because they constitute a natural extension of the notions and techniques
already developed. One will note as a consequence that the reading of
expos\'es~$\mathrm{VII}_{\mathrm{A}}$, $\mathrm{VII}_{\mathrm{B}}$, XV, XVI,
XVII is not necessary for the study of the rest of the Seminar, and in
particular for the part of this Seminar devoted to reductive group schemes.

\par\addvspace{0.55\baselineskip}
\noindent\textbf{6.}\label{Intro.sec.6}
From the theory of schemes, we shall use above all the general language of
schemes, set out in EGA~I, the notions of flat morphism, \'etale morphism,
smooth morphism, set out in SGA~1, I to~V, and finally the theory of
faithfully flat descent of SGA~1, VIII. We have as far as possible avoided
formulating unnecessary noetherian hypotheses, which has on the other hand
obliged us to replace the usual hypothesis ``of finite type'' by the
hypothesis ``of finite presentation''. For the notion of morphism of finite
presentation, the reader will consult EGA~IV, 1.4 and~1.6. The results of
SGA~1, I to~IV, stated most often in the noetherian context, will be
developed in the general case in EGA~IV%
{\renewcommand{\thefootnote}{(****)}\footnote{Since the writing of this
introduction, the four parts (\S\S~1 to~21) of EGA~IV have appeared.}},
where standard methods will likewise be developed in detail for reducing
certain types of statements (involving hypotheses of finite presentation) to
the noetherian case (EGA~IV, paragraphs~8, 9, 11). The reader who does not
wish to admit these results of EGA~IV may simplify certain statements or
their proof by supposing the base prescheme to be locally noetherian. He
does, however, expose himself to difficulties in the cases where the proof
proceeds by descent from~$\widehat{A}$ to~$A$, where~$\widehat{A}$ is the
completion of a noetherian local ring~$A$, for this method leads one to
introduce the ring (in general non-noetherian)
\(\widehat{A}\otimes_{A}\widehat{A}\).

\par\addvspace{0.55\baselineskip}
\noindent\textbf{7.}\label{Intro.sec.7}
References will be made according to the usual decimal system: the
reference~5.7.11 refers to the proposition (or lemma, definition, etc.)\ of
that name in the same expos\'e; in the reference~XVII~7.8 the Roman numeral
indicates the number of the expos\'e. We shall use the following sigla for
our standard references:

\medskip
\noindent
\begin{tabular}{@{}lp{0.72\textwidth}@{}}
Bible
  & $=$ S\'eminaire Chevalley ``Groupes de Lie alg\'ebriques'', 1956/58.\\[0.45em]
EGA~$X$, $x.y.z$
  & $=$ J.~Dieudonn\'e and A.~Grothendieck, \emph{\'El\'ements de G\'eom\'etrie
    Alg\'ebrique}, Chap.~$X$, statement~$x.y.z$ (or subsection~$x.y$, etc.).\\[0.45em]
SGA~$n$, $X$~$y.z$
  & $=$ S\'eminaire de G\'eom\'etrie Alg\'ebrique du Bois-Marie, year~$n$,
    statement~$y.z$ of expos\'e~$X$.\\[0.45em]
TDTE
  & $=$ A.~Grothendieck, \emph{Techniques de descente et Th\'eor\`emes
    d'existence en G\'eom\'etrie Alg\'ebrique}, expos\'es given in the
    Bourbaki Seminar between 1959 and 1962.
\end{tabular}

\medskip
{\renewcommand{\thefootnote}{(1)}%
\nde{Concerning the quotient by a reductive group, the situation has
evolved considerably since the writing of Note~(**) by A.~Grothendieck.
Indeed, for an affine algebraic group over an arbitrary field~$k$, the
``good'' notion, introduced by Mumford, is that of a ``geometrically
reductive'' group. (One says, on the other hand, that~$G$ is ``linearly
reductive'' if every rational representation of~$G$ is completely reducible,
but, as pointed out in Note~(**), this condition is too restrictive if
\(\operatorname{car}(k) > 0\)). According to the results of M.~Nagata and
W.~J.~Haboush (\cite{Na64}, \cite{NM64}, \cite{Ha75}), the geometrically
reductive \(k\)-groups are exactly the reductive \(k\)-groups in the sense
of the present Seminar. For all of this, see the later editions of Mumford's
book (\cite{MF82}). Moreover, the extension to the case of an arbitrary base
of the notion of ``geometric reductivity'', and of its consequences for
passage to the quotient, has been carried out by C.~S.~Seshadri
(\cite{Se77}), and additions, due to M.~Raynaud, are to be found in the
article~\cite{CTS79} of J.-L.~Colliot-Th\'el\`ene and J.-J.~Sansuc. Finally,
for more recent developments concerning passage to the quotient by an
algebraic group or, more generally, by a groupoid (\cf Expos\'e~V of the
present Seminar), let us mention the articles~\cite{Ko97} and~\cite{KM97} of
J.~Koll\'ar and of S.~Keel and S.~Mori.}}

\renewcommand{\bibname}{Bibliography}
\begin{thebibliography}{CTS79}

\bibitem[CTS79]{CTS79}
J.-L.~Colliot-Th\'el\`ene \& J.-J.~Sansuc, Fibr\'es quadratiques et
composantes connexes r\'eelles, \emph{Math.\ Ann.} \textbf{244} (1979),
105--134.

\bibitem[Ha75]{Ha75}
W.~J.~Haboush, Reductive groups are geometrically reductive, \emph{Ann.\ of
Math.} \textbf{102} (1975), \No 1, 67--83.

\bibitem[KM97]{KM97}
S.~Keel \& S.~Mori, Quotient by groupoids, \emph{Ann.\ of Math.}
\textbf{145} (1997), \No 1, 193--213.

\bibitem[Ko97]{Ko97}
J.~Koll\'ar, Quotient spaces modulo algebraic groups, \emph{Ann.\ of Math.}
\textbf{145} (1997), \No 1, 33--79.

\bibitem[MF82]{MF82}
D.~Mumford \& J.~Fogarty, \emph{Geometric invariant theory}, 2nd~ed.,
Springer-Verlag, 1982; (\resp 3rd~ed., with F.~Kirwan, 1994).

\bibitem[Na64]{Na64}
M.~Nagata, Invariants of a group in an affine ring, \emph{J.\ Math.\ Kyoto
Univ.} \textbf{3} (1964), \No 3, 369--377.

\bibitem[NM64]{NM64}
M.~Nagata \& T.~Miyata, Note on semi-reductive groups, \emph{J.\ Math.\ Kyoto
Univ.} \textbf{3} (1964), \No 3, 379--382.

\bibitem[Se77]{Se77}
C.~S.~Seshadri, Geometric reductivity over an arbitrary base, \emph{Adv.\
Math.} \textbf{26} (1977), \No 3, 225--274.

\end{thebibliography}
